\section{Multiplicity of the surface in the modular quotient}

The following argument uses the non-CM case of the pinned modular
realization, rather than merely the conclusion that the surface is
an isogeny factor. Let $A/L'$ be the abelian surface in that
construction. Its geometric endomorphism algebra is a quaternion
division algebra, so $A$ is absolutely simple. Let $f$ be the
primitive weight-two newform, let $E_f$ be its coefficient field,
let $\varepsilon$ be its nebentypus, and set
$d_f=[E_f:\mathbb Q]$ and $s=\operatorname{ord}(\varepsilon)$.
The scalar splitting construction gives, at one Tate prime,
\[
 V_\ell(A)\otimes\overline{\mathbb Q}_\ell
       \simeq r_{f,\ell}|_{G_{L'}}^{\oplus2}.
\]
The associated modular abelian variety $A_f$ has dimension $d_f$,
is an isogeny factor of $J_H=\operatorname{Jac}(X_1(N)/\ker\varepsilon)$
over $\mathbb Q$, and has Tate representation
\[
 V_\ell(A_f)\otimes\overline{\mathbb Q}_\ell
       \simeq\bigoplus_{\tau:E_f\hookrightarrow\overline{\mathbb Q}_\ell}
                         r_{f,\ell}^{\tau}.
\]

\begin{lemma}
Every coefficient conjugate $r_{f,\ell}^{\tau}|_{G_{L'}}$ is
isomorphic to $r_{f,\ell}|_{G_{L'}}$. Consequently $d_f$ is even
and
\[
                         (A_f)_{L'}\sim_{L'} A^{d_f/2}.
\]
\end{lemma}
\begin{proof}
At every good finite place $v$ of $L'$ away from $\ell$,
\[
 \operatorname{Tr}(r_{f,\ell}(\operatorname{Frob}_v))
       =\tfrac12\operatorname{Tr}
                    (V_\ell(A)(\operatorname{Frob}_v))\in\mathbb Q.
\]
Compatibility of the modular system identifies this trace with an
element of $E_f$, independently of the chosen coefficient embedding.
For a place of residue degree $f_v$ above a rational prime, it is
the corresponding power-trace polynomial in the usual modular
Frobenius coefficients. Hence every embedding $\tau$ fixes the
displayed trace. Chebotarev density and semisimplicity give the
asserted isomorphism of restricted representations.

It follows that the Tate modules of $(A_f)^2$ and $A^{d_f}$
become isomorphic over $\overline{\mathbb Q}_\ell$. They are
semisimple representations over $\mathbb Q_\ell$, so their
characters already identify them over $\mathbb Q_\ell$. Faltings'
isogeny theorem gives $(A_f)^2\sim_{L'}A^{d_f}$. Absolute
simplicity of $A$ and uniqueness of the simple factors in the
isogeny category show that $A_f\sim_{L'}A^r$ for an integer $r$.
Dimensions give $d_f=2r$.
\end{proof}

\begin{lemma}
The coefficient field $E_f$ contains $\mathbb Q(\zeta_s)$.
Consequently, if $s$ is a power of two, the multiplicity $r=d_f/2$
satisfies $r\ge s/4$, and also $r\ge1$.
\end{lemma}
\begin{proof}
For every prime $p\nmid N$ the Hecke relation gives
\[
               a_p(f)^2-a_{p^2}(f)=p\varepsilon(p).
\]
Thus $E_f$ contains every $\varepsilon(p)$ at such primes.
Primes in residue classes modulo the conductor supply all values
of the character, so $E_f$ contains the full group of $s$th roots
of unity. Therefore $d_f\ge\varphi(s)$. If $s\ge2$ is a power
of two, this gives $r\ge s/4$. For $s=1$, $r\ge1$ suffices.
\end{proof}

\section{Height averaging through surface quotients}

\begin{lemma}
Let $J/L'$ have dimension $g$, and suppose $A^r$ is an isogeny
factor of $J$ over $L'$, with $A$ an abelian surface. There is an
abelian surface $A_0/L'$ isogenous to $A$ such that
\[
 h_F(A_0)\le
          \frac{h_F(J)+(3/2)(g-2r)}{r}
       \le\frac{h_F(J)+3g/2}{r}.
\]
All heights are stable Faltings heights in the original normalization.
\end{lemma}
\begin{proof}
Let $B\subset J$ be the image of an $L'$-isogeny from $A^r$.
Let $B_i$ be the image of the first $i$ factors, with $B_0=0$.
Then each $C_i=B_i/B_{i-1}$ is an abelian surface over $L'$
isogenous to $A$. Repeated application of Rémond's exact-sequence
inequality gives
\[
                  \sum_{i=1}^r h_F(C_i)\le h_F(B).
\]
Another application gives
$h_F(B)+h_F(J/B)\le h_F(J)$.
Bost's lower bound in this normalization is
$h_F(J/B)>-(3/2)(g-2r)$, yielding the first displayed upper
bound for the average of the $h_F(C_i)$. One of them is no
larger than the average. Take it for $A_0$.
\end{proof}

Set
\[
 \mathcal H_r=\max\left(1,\frac{h_F(J)+3g/2}{r}\right),
 \qquad d'=[L':\mathbb Q].
\]
Applying the existing dimension-two Gaudron--Rémond isogeny
bound to $A_0\to A$ therefore gives
\[
 h_F(A)\le \mathcal H_r+
 8\left(32\log14+\log d'
                   +\log\max(1,\log d',\mathcal H_r)\right).
\]
No isogeny estimate in the dimension of $J$ or $A_f$ is used.

For the modular curve, its degree bound has the form
$d_H\le Ns\prod_{p\mid N}(1+1/p)$. The multiplicity estimate
$s/r\le4$ consequently removes the character-order factor from
$d_H/r$. Combined with the existing differential-evaluation height
bound, this gives an upper bound of order
\[
 R^2\prod_{p\mid R}(1+1/p)\,(1+\log R)
\]
for the non-CM surface, with effective constants determined by the
already explicit level and realization-field bounds.
For every fixed $\eta>0$, the elementary estimates for the finite
Euler product and logarithm then give a bound of the form
$C_\eta R^{2+\eta}$.

The small-height CM branch of the inherited modular-realization
input remains separate and already has a bound linear in $R$.

\paragraph{Primary source check.}
The pinned source, Version9 at commit
14e8f601fe0d9173d4a5f084aedc3e07066a7688, explicitly states
$\operatorname{End}^0(A_{\overline{\mathbb Q}})=B$ with $B$ a
quaternion division algebra and the identity
$\rho_\ell|_{G_{L'}}=V_\ell(A)$ with
$\rho_\ell=r_\ell\oplus r_\ell$ in its level-repair proof.
Rémond, \emph{Propriétés de la hauteur de Faltings}, Corollary1.2,
states $h_F(B)+h_F(J/B)\le h_F(J)$; the paragraph immediately
following Corollary1.3 states that the result also holds for stable
heights. These exact assertions were checked in the primary text
at https://www-fourier.univ-grenoble-alpes.fr/~remond/2476.pdf.
